\documentclass[a4paper,11pt]{article}
\usepackage{exam-style-842-unit}
\examinehead{842 信息化概论 · 第 05 单元练习册}{练习三 参考答案与解析}
\begin{document}

\unitanswerscover{842 信息化概论 · 单元练习册}{第 05 单元参考答案}{练习三（大数据与 Hadoop）的答案与解析；判断题均给改正表述，教材未展开的 4V／5V、YARN、Shuffle、总体架构与可视化均注明来源。}

\answerpart{一、单项选择题（10 题，每题 3 分，共 30 分）}

\q D。\begin{analysis}教材 4.2.2 指出狭义大数据具有 Volume（数据量）、Variety（种类）、Velocity（速度）3 个特征（第 53—54 页，图 4.5）。Value（价值）是通行的“4V”说法在大数据特征中增加的第四项，加上 Veracity（真实性）即通行的“5V”说法；这两项均非教材原文，故本册判断题、选择题以教材 3V 为准，4V／5V 作为补充了解。\end{analysis}

\q A。\begin{analysis}教材 4.2.2 给出大数据定义：大数据（Big Data 或 Megadata），或称巨量数据、海量数据、大资料，指的是所涉及的数据量规模巨大到无法通过人工在合理时间内被截取、管理、处理并整理，从而成为人类所能解读的形式的信息（第 53 页）。B 数据仓库是面向分析的集成化数据环境；C 数据库是结构化数据的存储管理软件；D 数据字典是描述数据的数据，均不符合。\end{analysis}

\q A。\begin{analysis}教材 4.2.1 指出：据统计，2013 年中国产生的数据总量超过 0.8ZB（相当于 8 亿 TB），相当于 2012 年数据总量的 2 倍，相当于 2009 年全球的数据总量（第 52 页）。B 8.5ZB 是预计 2020 年中国数据总量的下限；C、D 数量级与教材不符。\end{analysis}

\q A。\begin{analysis}教材 4.2.1 给出数据单位换算关系：计算机存储的最小基本单位是 Byte，按 1024（2 的 10 次方）进率计算，单位依次为 bit、Byte、KB、MB、GB、TB、PB、EB、ZB、YB、DB、NB；其中 1ZB = 1024EB（第 52 页）。B 1000、C 512、D 2048 均不符合 1024 进率。\end{analysis}

\q A。\begin{analysis}教材 4.2.3 指出：Hadoop 是根据 Google 公司发表的 MapReduce 和 Google 档案系统（即 GFS）的论文制作而成的（第 55 页）。B 是网络协议；C 是关系数据库理论；D 是人工智能技术，均与 Hadoop 的来源无关。\end{analysis}

\q A。\begin{analysis}HDFS 是一个主/从（Master/Slave）体系结构：HDFS 集群拥有一个 NameNode 和一些 DataNode，NameNode 管理文件系统的元数据，DataNode 存储实际的数据，客户端通过 NameNode 和 DataNode 的交互访问文件系统（教材 4.2.3，第 55—56 页）。C JobTracker 是旧版 MapReduce 的资源管理组件（教材未述）；D Hive 是数据仓库工具。\end{analysis}

\q C。\begin{analysis}教材 4.2.3 指出：HDFS 中的任意一份数据，理论上都至少存放 3 份以上的副本，任何一份副本的失效都不会影响数据访问（第 56 页）。这是 HDFS“可靠的数据存储”特点，A、B 副本过少起不到容错作用，D 并非教材表述。\end{analysis}

\q A。\begin{analysis}教材 4.2.3 列出 Hadoop 的知名子项目：Hive 是数据仓库工具，由 Facebook 贡献；Zookeeper 是分布式锁设施，也由 Facebook 贡献；Pig 是大数据分析平台；Avro 是数据序列化格式与传输工具；Ambari 是管理工具（第 57 页）。B HBase 是建立在 HDFS 之上的数据库；C Mahout 是机器学习算法库；D Storm 是实时计算框架，均不是“数据仓库工具”。\end{analysis}

\q A。\begin{analysis}教材 4.2.3 指出：相对于 Hadoop 的 MapReduce 在运行完工作后将中介数据存放到磁盘中，Spark 使用了存储器内运算技术，能在数据尚未写入硬盘时分析运算；内存运算速度比 Hadoop MapReduce 快约 100 倍，硬盘运算也快约 10 倍，适合机器学习算法（第 58—59 页）。B 错误：Spark 仍需搭配集群管理员与分布式存储系统；C 错误：Spark 可处理多种数据；D 错误：Spark 的 MLlib 正是机器学习库。\end{analysis}

\q A。\begin{analysis}教材 4.2.4 开篇即指出：大数据的发展并不是一片叫好之声，它同样面临着一些严重的问题，其中最严重的问题是隐私问题（第 60 页）。B 存储成本、C 带宽是可随技术进步缓解的工程问题；D 与大数据问题无关。\end{analysis}

\answerpart{二、多项选择题（4 题，每题 4 分，共 16 分）}

\q ABC。\begin{analysis}教材 4.2.1 把与企业经营相关的大数据来源划分为：物联网、Internet 查询数据、社会化数据、视频和音频 4 类（第 53 页）。A、B、C 分别对应前三类；D 单一结构化报表恰恰不是大数据的典型来源——大数据强调海量、非结构化、实时的数据，传统结构化报表属于较小规模、已加工的数据。\end{analysis}

\q ABC。\begin{analysis}教材 4.2.2 给出广义大数据定义：所谓大数据，是一个综合性概念，它包括因具备体积/种类/速度（Volume/Variety/Velocity，3V）特征而难以进行管理的数据，对这些数据进行存储、处理、分析的技术，以及能够通过分析这些数据获得实用意义和观点的人才和组织（第 54 页）。A、B、C 正确；D 错误，广义大数据明确包括非结构化数据，并非仅指关系数据库中的结构化数据。\end{analysis}

\q ABC。\begin{analysis}教材 9.6.13 指出：DM 与 KDD 现已成为人工智能应用的热门领域和研究方向，其涉及范围非常广泛，如企业数据、商业数据、科学实验数据、管理决策数据、Web 数据的挖掘和发现（第 159 页）。A、B、C 在列；D“仅限科学实验数据”缩小了范围，与教材表述不符。\end{analysis}

\q ABC。\begin{analysis}教材 4.2.3 列出的 Hadoop 子项目包括 Hadoop Common、Hive、Zookeeper、Avro、Pig、Ambari 等（第 57 页）。A Hive、B Pig、C Zookeeper 均在列；D 微信是腾讯的社交应用，属 SaaS，不是 Hadoop 子项目。\end{analysis}

\answerpart{三、判断题（6 题，每题 2 分，共 12 分）}

\q √。\begin{analysis}HDFS 是主/从体系结构，集群拥有一个 NameNode 和一些 DataNode：NameNode 管理文件系统的元数据，DataNode 存储实际的数据（教材 4.2.3，第 55—56 页）。\end{analysis}

\q ×。改正：MapReduce 计算时间通常以小时为单位，是事实上的批处理框架，适用于离线计算；实时计算可用 Storm 等流式计算框架。\begin{analysis}教材 4.2.3 指出：以 MapReduce 为基础的计算花费时间较长，通常以小时为单位，这决定了它只能用于离线计算环节，也就是说 MapReduce 是事实上的批处理框架；大数据实时计算是另一发展方向，Storm 是分布式的、容错的实时计算流式框架（第 57—58 页）。\end{analysis}

\q √。\begin{analysis}Spark 使用存储器内运算技术，能在数据尚未写入硬盘时分析运算；在存储器内运行程序的运算速度能做到比 Hadoop MapReduce 快约 100 倍，即便运行于硬盘也能快约 10 倍（教材 4.2.3，第 58—59 页）。\end{analysis}

\q ×。改正：隐私是大数据最严重的问题；个人数据可能被随意收集和滥用，涉及数据应当属于谁、谁有权利利用这些数据进行分析、这种利用是否应有限度 3 个问题。\begin{analysis}教材 4.2.4 用“必胜客客服”的段子说明“大数据时代没有隐私”，并归纳出上述 3 个问题（第 60—62 页），还列举了斯诺登事件与棱镜计划。\end{analysis}

\q √。\begin{analysis}教材 4.2.3 指出：Hadoop 是一款支持数据密集型分布式应用并以 Apache 2.0 许可协议发布的开源软件框架（第 55 页）；其优点之一就是低成本，因为 Hadoop 是开源的。\end{analysis}

\q ×。改正：Hive 是数据仓库工具，Zookeeper 是分布式锁设施（提供类似 Google Chubby 的功能）。\begin{analysis}教材 4.2.3 的子项目列表中，Hive 为“数据仓库工具，由 Facebook 贡献”，Zookeeper 为“分布式锁设施，提供类似 Google Chubby 的功能，由 Facebook 贡献”（第 57 页）。题干将二者作用互换。\end{analysis}

\answerpart{四、配对题（6 对，每对 1 分，共 6 分）}

\q （1）—B\quad （2）—D\quad （3）—A\quad （4）—C\quad （5）—E\quad （6）—F。
\begin{analysis}HDFS 是分布式文件系统，由 NameNode（元数据）与 DataNode（实际数据）组成主/从结构（B）；MapReduce 是分布式计算框架，合并映射与化简两种经典函数（D）；Hive 是数据仓库工具，由 Facebook 贡献（A）；Zookeeper 是分布式锁设施，提供类似 Google Chubby 的功能（C）；Spark 是采用存储器内运算的开源计算框架（E）；Mahout 是提供可扩展机器学习经典算法实现的 Hadoop 子项目（F）。（教材 4.2.3，第 55—59 页）\end{analysis}

\answerpart{五、简答题（3 题，每题 8 分，共 24 分）}

\q 参考答案要点：
（1）教材定义（2 分）：大数据指的是所涉及的数据量规模巨大到无法通过人工在合理时间内被截取、管理、处理并整理，从而成为人类所能解读的形式的信息（教材 4.2.2，第 53 页）。
（2）狭义定义（2 分）：所谓大数据，就是用现有的一般技术难以管理的大量数据的集合；狭义大数据只针对数据本身（第 53—54 页）。
（3）三个特征（3 分）：Volume（数据量）——用现有技术无法管理的数据量，目前基本指从几十 TB 到几 PB 的数量级，且随技术进步不断增加；Variety（种类）——既包括传统结构化数据，也包括网站日志、社交媒体文本、GPS 位置信息、传感器数据、图片和视频等非结构化数据；Velocity（速度）——数据产生和更新的频率，如便利店 24 小时的 POS 数据、网站点击流、高峰时每秒 7000 条的 Twitter 推文等。
（4）广义定义（1 分）：广义大数据是一个综合性概念，包括 3V 数据、存储处理分析这些数据的技术，以及能从中获得实用意义和观点的人才和组织（第 54 页）。
（5）可补充：通行的 4V 说法在 3V 基础上增加 Value（价值），5V 说法再增加 Veracity（真实性），但教材只列 3V，答题应以教材为准并说明补充来源。

\q 参考答案要点（教材 4.2.3，第 56—58 页）：
（1）要解决的问题（1 分）：HDFS 解决分布式存储问题，MapReduce 解决分布式计算问题；以 C++ 程序为主的传统分布式计算编程非常复杂、编写调试难度大，MapReduce 的出现极大缓解了这一问题。
（2）原理（3 分）：MapReduce 合并了两种经典函数——映射（Mapping）和化简（Reducing）。映射对集合里的每个目标应用同一个操作，即把大任务拆成可并行的小任务；化简遍历集合中的元素来返回一个综合的结果，即把并行结果汇总。教材用“数黑桃”为例：把牌分给所有玩家（Step1），每个玩家数自己手中的黑桃并汇报（Step2，Map），把所有玩家报出的数字加起来得到结论（Step3，Reduce）。
（3）作用（2 分）：Hadoop 的 MapReduce 是一个使用简易的分布式计算框架，基于它写出来的应用程序能够运行在由上千个商用机器组成的大型集群上，并以一种可靠容错的方式并行处理 TB 级的数据集；它可使用 Java 语言编程，降低了程序设计的门槛，规避了 C++ 容易出现内存问题的风险，获得广泛使用。
（4）执行流程（2 分，教材未展开 Shuffle，按通行资料补充）：输入数据被划分为分片 $\rightarrow$ Map 阶段各节点并行处理本地数据、输出键值对 $\rightarrow$ Shuffle 阶段按 key 进行分区、排序与分组，把相同 key 的数据送到同一 Reduce 节点 $\rightarrow$ Reduce 阶段汇总计算并输出结果。教材未使用“Shuffle”一词，只讲映射与化简两个阶段，答题时应以 Map、Reduce 为主线并说明中间的分组排序。
（5）为什么能提高速度（1 分）：任务被拆分为可并行的子任务，机器越多算得越快；多个节点并行处理、结果汇总，且具备容错与副本机制，适合批处理海量数据。

\q 参考答案要点（教材 4.2.3，第 55—60 页；YARN 为通行补充）：
（1）HDFS（2 分）：Hadoop 分布式文件系统，位于 Hadoop 底层，存储集群中所有存储节点上的文件；采用主/从结构，NameNode 管元数据、DataNode 存数据，数据多副本、可处理超大文件、流式访问、运行于廉价商用机器集群。它解决“数据存在哪里”的问题。
（2）YARN（1 分）：Hadoop 的资源管理与作业调度层，负责集群 CPU、内存等资源的统一管理与任务调度。教材仅在 Spark 的管理方式处提到“Hadoop YARN”，未系统介绍，本册按通行资料补充；答题时说明其“资源调度”定位即可。
（3）MapReduce（2 分）：分布式计算框架，用映射与化简完成批处理计算，可靠容错，是 Hadoop 最早的计算引擎；它解决“怎么算”的问题。
（4）HBase（1 分）：建立在 HDFS 之上的数据库，提供数据库的相关功能，适合海量数据的随机读写（教材 4.2.3 指出“在 HDFS 之上的 HBase/Hive 提供了数据库的相关功能”）。
（5）Hive（1 分）：建立在 HDFS 之上的数据仓库工具，由 Facebook 贡献，把数据组织成表以支持分析查询。
（6）Spark（1 分）：开源分布式计算框架，采用存储器内运算，速度比 MapReduce 快，适合机器学习等迭代计算；可运行在 YARN 等资源管理器之上，并可与 HDFS 等分布式存储搭配。
（7）关系：HDFS 是底层存储，YARN 负责资源调度，MapReduce 与 Spark 是计算引擎，HBase 与 Hive 在 HDFS 之上分别提供数据库与数据仓库功能；它们共同构成 Hadoop 生态。

\answerpart{六、综合题（1 题，共 12 分）}

\q 参考答案要点：
（1）来源归属与 3V（4 分）：
来源（2 分）：第一，田间传感器与农机作业数据属于物联网数据（机器连续测量并报告运行情况）；第二，气象与遥感影像属于专业研究机构与传感器观测数据，教材的来源四分类没有单列遥感，可归入物联网（传感器）或专业研究机构；第三，电商用户评论属于社会化数据，交易记录属于企业经营数据（也可归入 Internet 查询/交易数据）；第四，园区监控视频属于视频和音频。答题时说明归类理由即可（教材 4.2.1，第 52—53 页）。
3V（2 分）：Volume——数据量从 TB 级到 PB 级，远超传统数据库处理能力；Variety——结构化（交易记录）、半结构化（日志）、非结构化（评论文本、遥感影像、监控视频）并存，多源异构；Velocity——传感器数据秒级采集、交易与评论实时产生、监控视频连续流式产生，要求快速处理。
（2）总体架构（4 分）：按“采集—存储—处理—分析—应用”分层设计（教材第 4 章未系统给出这一分层，本小题按考试大纲与通行资料整理，组件作用依据教材 4.2.3）：
采集层：传感器网络、农机终端、业务系统日志、网络爬虫与数据接口，完成数据接入与预处理（清洗、去重、格式转换）。
存储层：用 HDFS 存储原始与历史海量数据，用 HBase 提供需要随机读写的数据库服务。
处理层：用 MapReduce 做批量离线计算，用 Spark 做内存计算与迭代计算，提高处理速度。
分析层：用 Hive 数据仓库组织数据并做查询统计，用 Mahout、MLlib 等做分类、聚类、推荐等数据挖掘。
应用层：面向精准种植、农机调度、电商选品与精准营销、灾害预警等业务提供决策支持与可视化展示。
（3）问题、对策与呈现（4 分）：
隐私（1.5 分）：存在数据应当属于谁、谁有权分析、利用是否应有限度 3 个问题，个人数据可能被随意收集和滥用（4.2.4，第 60—62 页）。对策：取得用户授权、对敏感数据脱敏与加密、最小权限访问、审计日志、遵守相关法律法规。
数据质量（1 分，教材未系统展开，按通行整理）：多源异构数据存在缺失、噪声、不一致、口径不统一等问题；对策是建立数据标准与元数据管理、数据清洗与校验、质量监控。
人才（0.5 分）：广义大数据强调需要能通过分析获得实用意义和观点的人才和组织；应培养既懂农业业务又懂数据技术的复合型人才，并建立数据治理组织。
可视化呈现（1 分）：可用折线图、柱状图、散点图、热力图、专题地图（结合 GIS）、仪表盘与数据大屏等方式，把分析结果直观呈现给管理者；教材 4.2.5 也提到百度旅游大数据平台对旅游业提供数据可视化和 API 接口服务（第 63 页）。

\end{document}
